\NativeSlot{40}{61}{%
\ExamGroup{三、计算。{\normalfont\kaishu\fontsize{10.8bp}{14bp}\selectfont （共30分）}}

% p107_s01

\ExamQuestion{1.}{5}%
直接写得数。（8分）\par
\noindent \NativeOralCell{3.7+4.3=}\NativeOralCell{8.36-5.3=}\NativeOralCell{0.49\times100=}\NativeOralCell{280\div5\div2=}\par
\noindent \NativeOralCell{9\div100=}\NativeOralCell{0.2\times1000=}\NativeOralCell{6.4+4.6=}\NativeOralCell{35\times25\times4=}\par
}
\NativeSlot{142}{84}{%
\par\noindent 2.\quad 列竖式计算，带☆的要验算。（10分）\par
\par\noindent%
\ExamTile{.323}{%
% p107_s02

\(16.8+2.27=\hspace{1.8em}\)\par

\ExamWork{7}
}\hfill%
\ExamTile{.323}{%
% p107_s03

☆\(12-4.5=\hspace{1.8em}\)\par

\ExamWork{7}
}\hfill%
\ExamTile{.323}{%
% p107_s04

\(10.02-1.9=\hspace{1.8em}\)\par

\ExamWork{5}
}\par
}
\NativeSlot{283}{75}{%
\par\noindent 3.\quad 计算下面各题，能简算的要简算。（12分）\par
\par\noindent%
\ExamTile{.485}{%
% p107_s05

\(180+19\times(106+74)\)\par

\ExamWork{3}
}\hfill%
\ExamTile{.485}{%
% p107_s06

\(40.65-(7.38+10.65)\)\par

\ExamWork{3}
}\par
}
\NativeSlot{424}{70}{%
\par\noindent%
\ExamTile{.485}{%
% p107_s07

\(25\times201-25\)\par

\ExamWork{3}
}\hfill%
\ExamTile{.485}{%
% p107_s08

\(840\div[15\times(81-77)]\)\par

\ExamWork{3}
}\par
}
\NativeSlot{542}{152}{%
\ExamGroup{四、操作。{\normalfont\kaishu\fontsize{10.8bp}{14bp}\selectfont （共6分）}}

\par\noindent%
\ExamTile{.54}{%
% p107_s09

\noindent 下图是一块平行四边形露营地，奇思在这块露营地的小路上跑步，淘气在草坪上跑步。\par （1）淘气跑的路线与小路平行，请你画一条直线表示淘气跑步的路线。\par

% p107_s10

\ExamQuestion{（2）}{5}%
笑笑想走到奇思所在的这条小路上，请你画出她去这条小路的最短路线。\par
}\hfill%
\ExamTile{.435}{%
\BookFigure{figures/p107_s09-01.tex}
}\par
}
